\chapter{Politica di sicurezza operativa}
\label{chap:sicurezza}

\section{Principi guida}
\label{sec:principi}

La politica di sicurezza si articola su cinque principi, ciascuno con realizzazione
concreta nel codice:

\begin{enumerate}[leftmargin=*]
    \item Principio di \emph{least privilege}
    \item \emph{Action boundaries} espliciti
    \item \emph{Idempotenza} delle mutazioni agentiche
    \item \emph{Auditability} completa
    \item \emph{Escalation \& rollback} strutturati
\end{enumerate}

\section{Principio di least privilege}
\label{sec:least_privilege}

\begin{itemize}[leftmargin=*]
    \item Il \textbf{MCP server} espone esclusivamente tool read-only sui metadati
          (folder tree, document inventory). Non esiste nessun tool di scrittura,
          eliminazione o spostamento esposto via MCP.
    \item L'\textbf{agente} non possiede credenziali JWT per il backend pubblico:
          ogni mutazione passa dall'agent worker attraverso operazioni bounded
          e tipizzate (\texttt{auto\_filed}, \texttt{sent\_to\_inbox}, ecc.).
    \item Le \textbf{credenziali email IMAP} sono cifrate a riposo con
          \textsc{Fernet} (AES-128-CBC con HMAC-SHA256). La chiave di cifratura
          \`e iniettata esclusivamente come variabile d'ambiente
          (\texttt{EMAIL\_ENCRYPTION\_KEY}) e non \`e mai persistita nei log.
    \item I \textbf{bucket S3} sono logicamente isolati per tenant tramite prefisso
          di path; il gateway S3 di SeaweedFS \`e configurato con un'identit\`a
          unica e ACL coerente (\texttt{s3.json}).
\end{itemize}

\begin{lstlisting}[language=Python, caption={Cifratura delle credenziali email (\texttt{backend/app/crypto.py}).}]
from cryptography.fernet import Fernet, InvalidToken

def encrypt_credentials(plain_text: str) -> str:
    """AES-128-CBC + HMAC-SHA256 (Fernet)."""
    f = Fernet(settings.email_encryption_key.encode())
    return f.encrypt(plain_text.encode()).decode()

def decrypt_credentials(encrypted_text: str) -> str:
    try:
        f = Fernet(settings.email_encryption_key.encode())
        return f.decrypt(encrypted_text.encode()).decode()
    except InvalidToken:
        raise ValueError("Credenziali non decifrabili: chiave errata o dati corrotti")
\end{lstlisting}

\section{Action boundaries}
\label{sec:action_boundaries}

Le decisioni dell'agente producono esattamente due esiti:

\begin{description}[leftmargin=*]
    \item[\texttt{auto\_filed}] solo se \texttt{confidence} $\geq 0.85$ \emph{e}
          \texttt{folder\_id} \`e un UUID Python valido \emph{e} quella chiave
          \`e una FK esistente nel DB del tenant. La validazione \emph{ex-post}
          dell'UUID copre il caso patologico in cui il modello produca un UUID
          sintatticamente valido ma semanticamente inesistente (hallucination).

    \item[\texttt{in\_inbox}] in tutti gli altri casi. Il documento resta accessibile
          nell'\emph{Inbox view} con tre azioni reversibili: \textbf{Accept},
          \textbf{Move} (a cartella diversa da quella proposta), \textbf{Reject}.
\end{description}

L'agente \textbf{non pu\`o}:
eliminare o spostare documenti gi\`a archiviati, modificare permessi, accedere a dati
di altri tenant, modificare account email o ruoli utenti. Tutte queste operazioni
richiedono autenticazione JWT umana e sono soggette al modello dei permessi del backend.

\section{Autenticazione e autorizzazione}
\label{sec:auth}

\begin{lstlisting}[language=Python, caption={Generazione e validazione JWT (\texttt{backend/app/auth.py}).}]
def create_access_token(user_id: str,
                        expires_delta: timedelta = None) -> str:
    expire = datetime.utcnow() + (expires_delta or timedelta(minutes=30))
    payload = {"sub": user_id, "exp": expire}
    return jwt.encode(payload, settings.secret_key, algorithm="HS256")

def get_current_user(request: Request,
                     db: Session = Depends(get_db)) -> User:
    # Estrae "Bearer {token}" dall'header Authorization
    # Decodifica JWT (verifica firma e scadenza)
    # Verifica che l'utente esista e sia attivo
    ...
\end{lstlisting}

Le password sono hash-ate con \textbf{bcrypt} tramite \texttt{passlib}:

\begin{lstlisting}[language=Python]
pwd_context = CryptContext(schemes=["bcrypt"], deprecated="auto")
pwd_context.hash(plain_password)    # all'inserimento
pwd_context.verify(plain, hashed)  # al login
\end{lstlisting}

\section{Debounce delle scritture di attivit\`a (\texttt{LastActiveMiddleware})}
\label{sec:middleware}

Il middleware \texttt{LastActiveMiddleware} aggiorna il campo
\texttt{users.last\_active\_at} al pi\`u una volta ogni 5~minuti per utente,
evitando una scrittura su ogni singola API call:

\begin{lstlisting}[language=Python, caption={Logica di debounce (\texttt{backend/app/middleware.py}).}]
_UPDATE_INTERVAL = timedelta(minutes=5)
_last_updated: dict[str, datetime] = {}   # cache in-memory per processo

class LastActiveMiddleware(BaseHTTPMiddleware):
    async def dispatch(self, request: Request, call_next) -> Response:
        response = await call_next(request)
        # ... estrae user_id dal JWT senza re-autenticare ...
        now = datetime.now(timezone.utc)
        last = _last_updated.get(user_id)
        if last and (now - last) < _UPDATE_INTERVAL:
            return response   # skip: aggiornato di recente
        # aggiorna DB e cache
        ...
        return response
\end{lstlisting}

\section{Isolamento multi-tenant}
\label{sec:multitenancy_sec}

L'isolamento \`e garantito a tre livelli:

\begin{enumerate}[leftmargin=*]
    \item \textbf{S3}: gli oggetti sono archiviati sotto il prefisso
          \texttt{\{domain-with-dashes\}/\{document-uuid\}}. Un accesso diretto
          a un path di altro tenant \`e prevenuto dalla configurazione del gateway S3.
    \item \textbf{Database}: ogni query sul backend filtra per \texttt{owner\_id}
          dell'utente autenticato. L'accesso cross-owner avviene solo tramite record
          esplicito nella tabella \texttt{permissions}.
    \item \textbf{Applicativo}: la condivisione tra utenti di domini diversi \`e vietata:
\begin{lstlisting}[language=Python]
user_domain = current_user.email.split("@")[1]
target_domain = target_user.email.split("@")[1]
if user_domain != target_domain:
    raise HTTPException(403, "Condivisione cross-domain non consentita")
\end{lstlisting}
\end{enumerate}

\section{Validazione input e prevenzione injection}
\label{sec:validation}

\begin{itemize}[leftmargin=*]
    \item \textbf{Pydantic v2}: valida automaticamente tutti gli schema di request.
          La violazione produce \texttt{400 Bad Request} con handler centralizzato in
          \texttt{backend/app/main.py}.
    \item \textbf{SQLAlchemy ORM}: tutte le query usano parametri named binding;
          nessuna concatenazione di SQL raw \`e presente nel codice applicativo.
    \item \textbf{Path traversal}: i nomi dei file sono sanitizzati con
          \texttt{os.path.basename(file.filename)} prima dell'archiviazione.
    \item \textbf{XSS}: React effettua l'escaping automatico di tutte le
          interpolazioni JSX; nessun uso di \texttt{dangerouslySetInnerHTML}.
    \item \textbf{Hallucination validation}: il campo \texttt{folder\_id} prodotto
          dall'LLM viene validato come UUID Python e come FK esistente nel DB prima
          di qualsiasi mutazione.
\end{itemize}

\section{Mappa azioni: deterministico, agentico, human-in-the-loop}
\label{sec:action_map}

\begin{table}[ht]
\centering
\small
\begin{tabularx}{\linewidth}{l X l}
\toprule
\textbf{Operazione} & \textbf{Chi decide / esegue} & \textbf{Reversibile?} \\
\midrule
Sincronizzazione IMAP               & Deterministica (poller)                     & s\`i \\
Estrazione allegati                 & Deterministica (worker)                      & s\`i \\
Classificazione semantica           & Agentica (LLM, $\theta = 0.85$)              & s\`i \\
Auto-filing (conf $\geq \theta$)    & Agentica $\rightarrow$ scrittura deterministica & s\`i (revisione utente) \\
Invio in inbox (conf $< \theta$)    & Deterministica (post-LLM)                    & --- (advisory) \\
Accettazione proposta               & Human-in-the-loop                            & s\`i (move/reject) \\
Spostamento manuale                 & Human-in-the-loop                            & s\`i \\
Cancellazione documento             & Human-in-the-loop (owner)                    & no (cascade) \\
Cancellazione cartella              & Human-in-the-loop (owner)                    & no (cascade) \\
Disattivazione utente               & Human-in-the-loop (admin)                    & s\`i \\
Cambio ruolo utente                 & Human-in-the-loop (admin)                    & s\`i \\
Revoca permission                   & Human-in-the-loop (owner)                    & s\`i \\
\bottomrule
\end{tabularx}
\caption{Mappa azioni: deterministico, agentico, human-in-the-loop.}
\label{tab:action_map}
\end{table}

\section{Anti-loop e guardrail economici}
\label{sec:guardrail}

Il worker \textbf{non implementa retry automatici illimitati}: in caso di eccezione
durante \texttt{\_process\_job}, il job transisce a \texttt{failed} con
\texttt{error\_message} e non viene rielaborato automaticamente. Questo elimina il rischio
di \emph{thrash} su job che producono sistematicamente errori (e.g.\ file corrotti,
API LLM down).

I cost guardrail sono realizzati da:
\begin{enumerate}[leftmargin=*]
    \item \textbf{Dimensione capped delle anteprime testuali} (Sezione~\ref{sec:eml_parser}):
          limita i token in input indipendentemente dalla dimensione del documento.
    \item \textbf{Batch size limitato}: il worker preleva al pi\`u 10~job per iterazione
          (\texttt{.limit(10)}).
    \item \textbf{Poll interval}: 30~secondi tra un ciclo e l'altro.
    \item \textbf{Structured output}: \texttt{instructor} forza la conformit\`a dello
          schema senza retry indefiniti.
\end{enumerate}
